\label{sec:limites}
\begin{definition} \textbf{L\'imite.} \label{def:limite_normas} 
Sean $\campoVec{F}:A \subseteq \Rn{n} \to \Rn{m}$ una funci\'on, $L \in \Rn{m}$ y $\vect{x}_0$ un punto de acumulaci\'on de $A$. Decimos que \emph{el l\'imite de $\campoVec{F}$ para $\vect{x}$ tendiendo a $\vect{x}_0$ es $L$} o que \emph{$\campoVec{F}$ tiende a $L$ cuando $\vect{x}$ tiende a $\vect{x}_0$} si 
\[ 
\forall \varepsilon > 0, \, \exists \, \delta > 0 \quad \text{tal que} \quad 0 < \lVert \vect{x} - \vect{x}_0 \rVert < \delta \land \vect{x} \in A \implies \lVert \campoVec{F}(\vect{x}) - L \rVert < \varepsilon. 
\] 
En este caso, utilizamos la notaci\'on 
\[ 
\campoVec{F}(\vect{x}) \xrightarrow[\vect{x} \to \vect{x}_0]{} L \quad \text{o} \quad \lim_{\vect{x} \to \vect{x}_0} \campoVec{F}(\vect{x}) = L. 
\] 
\end{definition}



\begin{obs} \label{obs:entornos_distancias}
A veces es útil expresar la misma definición en t\'erminos de entornos (bolas), lo cual equivale a utilizar funciones de distancia:
\[ 
\forall \varepsilon > 0, \, \exists \, \delta > 0 \quad \text{tal que} \quad \vect{x} \in B_{\delta}^*(\vect{x}_0) \cap A \implies \campoVec{F}(\vect{x}) \in B_{\varepsilon}(L). 
\] 
Recordemos que la pertenencia a la bola reducida $\vect{x} \in B_{\delta}^*(\vect{x}_0)$ equivale matem\'aticamente a la desigualdad de distancias $0 < \dis_n(\vect{x},\vect{x}_0) < \delta$.
\end{obs}

\begin{obs} \label{obs:campo_escalar} 
Si la funci\'on es un \emph{campo escalar} (es decir, $m = 1$), se suele denotar a la funci\'on con la letra min\'uscula $f$ y a su l\'imite con la letra $l \in \mathbb{R}$. En este caso, la condici\'on sobre el codominio $\lVert \campoVec{F}(\vect{x}) - L \rVert < \varepsilon$ se reduce al valor absoluto tradicional, escribi\'endose la definici\'on como:
\[ 
\forall \varepsilon > 0, \, \exists \, \delta > 0 \quad \text{tal que} \quad 0 < \lVert \vect{x} - \vect{x}_0 \rVert < \delta \land \vect{x} \in A \implies \abs{f(\vect{x}) - l} < \varepsilon. 
\] 
\end{obs}


\begin{theorem}\textbf{Unicidad del l\'imite.} \label{teo:unicidad_limite} Sean $\mathbf{F}:A  \subseteq  \Rn{n} \to \Rn{m}$ una funci\'on, $\mathbf{x}_0$ un punto de acumulaci\'on de $A$ y $\mathbf{L}_1, \mathbf{L}_2 \in \Rn{m}$ tales que 
\[
 \mathbf{F}(\mathbf{x}) \xrightarrow[\mathbf{x} \to \mathbf{x}_0]{} \mathbf{L}_1 \quad \wedge \quad \mathbf{F}(\mathbf{x}) \xrightarrow[\mathbf{x} \to \mathbf{x}_0]{} \mathbf{L}_2,
\]
entonces  $\mathbf{L}_1 = \mathbf{L}_2$.
\end{theorem}


\begin{proof}  Como 
\[
 \mathbf{F}(\mathbf{x}) \xrightarrow[\mathbf{x} \to \mathbf{x}_0]{} \mathbf{L}_1,
\]
dado $\varepsilon > 0$ existe alg\'un $\delta_1 > 0$ tal que cada punto $\mathbf{x} \in B_{\delta_1}^*(\mathbf{x}_0) \cap A$ se aplica a trav\'es de $\mathbf{F}$ en alg\'un punto de la bola de centro $\mathbf{L}_1$ y radio $\varepsilon$.

\def\Ro{2.0}
\def\Rb{1.0}
\def\Ra{3.0}
\def\sep{1.0*\Ro}

\input{figs/fig1}

De la misma manera, como 
\[
 \mathbf{F}(\mathbf{\mathbf{x}}) \xrightarrow[\mathbf{\mathbf{x}} \to \mathbf{\mathbf{x}}_0]{} \mathbf{L}_2,
\]
dado $\varepsilon > 0$ existe alg\'un $\delta_2 > 0$ tal que cada punto $\mathbf{\mathbf{x}} \in B_{\delta_2}^*(\mathbf{\mathbf{x}}_0) \cap A$ se aplica a trav\'es de $\mathbf{F}$ en alg\'un punto de la bola de centro $\mathbf{L}_2$ y radio $\varepsilon$.

\input{figs/fig2}

Si suponemos que $\mathbf{L}_1 \ne \mathbf{L}_2$, por propiedad de la distancia es $\dis(\mathbf{L}_1,\mathbf{L}_2) > 0$ y podemos elegir 
\[
 \varepsilon = \frac{\dis{(\mathbf{L}_1,\mathbf{L}_2)}}{2}.
\]
Para este valor de $\varepsilon$ podemos encontrar un punto $\mathbf{\mathbf{x}}$ que se aplica por $\mathbf{F}$ en alg\'un punto de la bola de centro $\mathbf{L}_1$ y radio $\varepsilon$ \textbf{y tambi\'en} se aplica por $\mathbf{F}$ en alg\'un punto de la bola de centro $\mathbf{L}_2$ y radio $\varepsilon$, lo cual es absurdo, pues estas dos bolas son disjuntas.

\input{figs/fig3}

Ahora s\'i, veamos la demostraci\'on formal del teorema.  Supongamos, por el absurdo, que $\mathbf{L}_1 \ne \mathbf{L}_2$, y sea $\varepsilon = \dis(\mathbf{L}_1,\mathbf{L}_2)/2 > 0.$  Como 
 \[
  \mathbf{F}(\mathbf{x}) \xrightarrow[\mathbf{x} \to \mathbf{x}_0]{} \mathbf{L}_1,
 \]
podemos elegir $\delta_1 > 0$ tal que $\mathbf{x} \in B_{\delta_1}^*(\mathbf{x}_0) \cap A \then \mathbf{F}(\mathbf{x}) \in B_{\varepsilon}(\mathbf{L}_1)$. Adem\'as, como 
\[
 \mathbf{F}(\mathbf{x}) \xrightarrow[\mathbf{x} \to \mathbf{x}_0]{} \mathbf{L}_2,
\]
 podemos elegir $\delta_2 > 0$ tal que $\mathbf{x} \in B_{\delta_2}^*(\mathbf{x}_0) \cap A \then \mathbf{F}(\mathbf{x}) \in B_{\varepsilon}(\mathbf{L}_2)$.  Notemos que, como sugiere la figura, 
 \[
  B_{\varepsilon}(\mathbf{L}_1) \cap B_{\varepsilon}(\mathbf{L}_2) = \varnothing.
 \]
 En efecto, de no ser as\'i, sea $z \in B_{\varepsilon}(\mathbf{L}_1) \cap B_{\varepsilon}(\mathbf{L}_2)$. Tenemos que:
 \[
  2 \varepsilon = \dis(\mathbf{L}_1,\mathbf{L}_2) \le \dis(\mathbf{L}_1,z) + \dis(z,\mathbf{L}_2)
  < \varepsilon + \varepsilon = 2 \varepsilon \then \varepsilon < \varepsilon. \text{ Absurdo.}
 \]

 Sea $\delta = \min\{ \delta_1, \delta_2 \}$. \\
 Como $B_{\delta}^*(\mathbf{x}_0)  \subseteq  B_{\delta_1}^*(\mathbf{x}_0)$ y $B_{\delta}^*(\mathbf{x}_0)  \subseteq  B_{\delta_2}^*(\mathbf{x}_0)$, si tomamos alg\'un $\mathbf{x} \in B_{\delta}^*(\mathbf{x}_0) \cap A$\footnote{Notemos que $B_{\delta}^*(\mathbf{x}_0) \cap A \ne \varnothing$ porque $\mathbf{x}_0$ es un punto de acumulaci\'on de $A$.}, vale que $f(\mathbf{x}) \in B_{\varepsilon}(\mathbf{L}_1)$ y tambi\'en $f(\mathbf{x}) \in B_{\varepsilon}(\mathbf{L}_2)$, por lo que $f(\mathbf{x}) \in B_{\varepsilon}(\mathbf{L}_1) \cap B_{\varepsilon}(\mathbf{L}_2)$. Absurdo, pues $B_{\varepsilon}(\mathbf{L}_1) \cap B_{\varepsilon}(\mathbf{L}_2) = \varnothing$. El absurdo provino de suponer que $\mathbf{L}_1 \ne \mathbf{L}_2$, por lo cual debe ser $\mathbf{L}_1 = \mathbf{L}_2$.
\end{proof}


\newpage



\begin{propertie}\textbf{\'Algebra de l\'imites.} \label{prop:alg_lim} 
Sean $\campoVec{F},\campoVec{G}:A \subseteq \Rn{n} \to \Rn{m}$ funciones y $\vect{x}_0$ un punto de acumulaci\'on de $A$. Supongamos que 
\begin{align*} 
&\lim_{\vect{x} \to \vect{x}_0} \campoVec{F}(\vect{x}) = L_1 \in \Rn{m}    \quad \text{ y } \\ 
&\lim_{\vect{x} \to \vect{x}_0} \campoVec{G}(\vect{x}) = L_2 \in \Rn{m},
\end{align*} 
entonces: 
\begin{enumerate} 
\item \textbf{Linealidad} 
\begin{enumerate} 
\item $\lim_{\vect{x} \to \vect{x}_0} (\campoVec{F} + \campoVec{G})(\vect{x}) = L_1 + L_2$ 
\item $\lim_{\vect{x} \to \vect{x}_0} (\alpha \campoVec{F})(\vect{x}) = \alpha L_1 \quad \forall \alpha \in \mathbb{R}$ 
\end{enumerate} 

\item $\lim_{\vect{x} \to \vect{x}_0} (\campoVec{F} \cdot \campoVec{G}) (\vect{x}) = L_1 \cdot L_2$\footnote{Si $m \ge 2$, ``$\cdot$'' representa el producto escalar en $\Rn{m}$.} 

\item Si $m = 1$ (\textit{i.e.}, si las funciones son campos escalares denotados por $f$ y $g$, con l\'imites $l_1$ y $l_2 \ne 0$), y $g$ es no nula en alg\'un entorno de $\vect{x}_0$, entonces: 
\[ 
\lim_{\vect{x} \to \vect{x}_0} \frac{f(\vect{x})}{g(\vect{x})} = \frac{l_1}{l_2}. 
\] 
\end{enumerate} 
\end{propertie}


\begin{tarea} Calcular $$ \lim_{(x,y) \to (0,0)}  (f(x) + g(x)),$$ siendo
 $$f(x,y) = 
     \begin{dcases}
         1  &  \text{ si } x < 0  \\
        -1                 & \text{ si }  x>0
     \end{dcases} \:\:\:  \text{ y }    \:\:\:  
  g(x,y) = 
     \begin{dcases}
         1  &  \text{ si } x > 0  \\
        -1                 & \text{ si }  x<0
     \end{dcases}$$   ¿Qu\'e puede decir sobre el uso de \'algebra de l\'imites ?
\end{tarea}

\begin{tarea} 
Calcular los siguientes l\'imites utilizando el \'algebra de l\'imites: 
\[ 
\lim_{(x,y) \to (0,0)} \frac{x^2 + 1 + y}{(x-1)^2 + y^2} 
\] 
\[ 
\lim_{(x,y) \to (0,0)} \frac{e^x \ln(\cos y)}{x^2 + (y+1)^2} 
\] 
\end{tarea}


\begin{obs} \label{obs:desigualdades_utiles} 
Para resolver l\'imites por definici\'on ser\'a \'util tener presente ciertas desigualdades elementales. Si $z \in \mathbb{R}$ y $(x,y) \in \mathbb{R}^{2} \setminus \{ (0,0) \}$, entonces: 
\begin{enumerate} 
\item $0 \le x^2 \le x^2+y^2 \iff 0 \le \frac{x^2}{x^2+y^2} \le 1 \iff 0 \le \frac{\abs{x}^2}{\lVert (x,y) \rVert^2} \le 1$ 
\item Tomando ra\'iz cuadrada en el \'item anterior se obtiene: \; $0 < \frac{\abs{x}}{\lVert (x,y) \rVert} \le 1 \iff \abs{x} \le \lVert (x,y) \rVert$.
\item $\abs{\sen z} \le 1 \quad \land \quad \abs{\cos z} \le 1  \quad  \quad \forall z \in \mathbb{R}$. 
\item $\abs{\sen z} \le \abs{z} \quad \forall z \in \mathbb{R} $. 
\end{enumerate} 
\end{obs}


\begin{example} 
Probar, usando la definici\'on de l\'imite, que 
\[ 
\lim_{(x,y) \to (0,0)} \frac{x^2 y}{x^2 + y^2} = 0. 
\] 
\end{example} 

\textbf{Soluci\'on.} Para el campo escalar $f: A=\mathbb{R}^{2} \setminus \{(0,0)\} \to \mathbb{R}$ definido por $f(x,y)=\frac{x^2 y}{x^2 + y^2}$ y su l\'imite $l=0$, debemos probar que 
\[ 
\forall \varepsilon > 0, \, \exists \, \delta > 0 \quad \text{tal que} \quad (x,y) \in A \land 0 < \lVert (x,y) \rVert < \delta \implies \abs{ \frac{x^2y}{x^2+y^2} - 0 } < \varepsilon. 
\] 
Partimos de la expresi\'on en valor absoluto:
\[ 
\abs{ \frac{x^2y}{x^2+y^2} - 0 } = \abs{ \frac{x^2y}{x^2+y^2} } = \frac{x^2 \abs{y}}{x^2+y^2} = \left( \frac{x^2}{x^2+y^2} \right) \abs{y}.
\] 
 aplicando los \'items 1 y 2 de la Observaci\'on \eqref{obs:desigualdades_utiles}, concluimos
\[ 
\left( \frac{x^2}{x^2+y^2} \right) \abs{y} \le 1 \cdot \abs{y} = \abs{y} \le \lVert (x,y) \rVert < \delta,
\]
luego, basta con elegir $\delta = \varepsilon$ para completar la demostraci\'on.

\begin{example} 
Probar, usando la definici\'on de l\'imite, que 
\[ 
\lim_{(x,y) \to (0,0)} \frac{x^2 \sen y}{x^2 + y^2} = 0. 
\] 
\end{example} 

\textbf{Soluci\'on.} Para el campo escalar $f: A=\mathbb{R}^{2} \setminus \{(0,0)\} \to \mathbb{R}$ definido por $f(x,y)=\frac{x^2 \sen y}{x^2 + y^2}$ y su l\'imite $l=0$, debemos probar que 
\[ 
\forall \varepsilon > 0, \, \exists \, \delta > 0 \quad \text{tal que} \quad (x,y) \in A \land 0 < \lVert (x,y) \rVert < \delta \implies \abs{ \frac{x^2 \sen y}{x^2+y^2} - 0 } < \varepsilon. 
\] 
Partimos de la expresi\'on en valor absoluto:
\[ 
\abs{ \frac{x^2 \sen y}{x^2+y^2} - 0 } = \abs{ \frac{x^2 \sen y}{x^2+y^2} } = \frac{x^2 \abs{\sen y}}{x^2+y^2}.
\] 
Utilizando el \'item 4 de la Observaci\'on \eqref{obs:desigualdades_utiles}, sabemos que $\abs{\sen y} \le \abs{y}$. Reemplazando obtenemos:
\[ 
\frac{x^2 \abs{\sen y}}{x^2+y^2} \le \frac{x^2 \abs{y}}{x^2+y^2} = \left( \frac{x^2}{x^2+y^2} \right) \abs{y}.
\] 
Nuevamente, aplicando los \'items 1 y 2 de la Observaci\'on \eqref{obs:desigualdades_utiles}, concluimos:
\[ 
 \left( \frac{x^2}{x^2+y^2} \right) \abs{y} \le  \lVert (x,y) \rVert < \delta.
\] 
Luego, basta con elegir $\delta = \varepsilon$.





%\newpage

\begin{propertie} \label{prop:lim_x_comp} 
Sea $\campoVec{F}: A \subseteq \Rn{n} \to \Rn{m}$ un campo vectorial y $\vect{x}_0$ un punto de acumulaci\'on de $A$. Sean $f_j: A \subseteq \Rn{n} \to \mathbb{R} \;\; (j = 1,2, \dots, m)$ campos escalares tales que $\campoVec{F}(\vect{x}) = \left( f_1(\vect{x}), f_2(\vect{x}), \dots, f_m(\vect{x}) \right), \; \forall \vect{x} \in A$\footnote{Los campos escalares $f_j$ quedan determinados un\'ivocamente por $\campoVec{F}$ de la siguiente manera: 
\begin{gather*} 
f_j : A \subseteq \Rn{n} \to \mathbb{R}; \\ 
f_j(\vect{x}) = \campoVec{F}(\vect{x}) \cdot \vect{e}_j, 
\end{gather*} 
donde $\vect{e}_j$ es el $j$-\'esimo vector can\'onico de $\Rn{m}$.}.\\ 
Sea $L = \left( l_1, l_2, \dots, l_m \right) \in \Rn{m}$, con $l_j \in \mathbb{R} \; \forall j = 1,2, \dots, m$. Entonces:
\[ 
\lim_{\vect{x} \to \vect{x}_0} \campoVec{F}(\vect{x}) = L \iff \lim_{\vect{x} \to \vect{x}_0} f_j(\vect{x}) = l_j \quad \forall j = 1,2, \dots, m. 
\] 
\end{propertie}

\begin{tarea}Calcular
 \[
  \lim_{(x,y) \to (0,0)}   \Big( \frac{x^2 \sin y}{x^2 + y^2},  \frac{x^2 +1 +  y}{(x-1)^2 +  y^2}, \frac{x^2  y}{x^2 + y^2}   \Big).
 \]
\end{tarea}

  \begin{propertie} \label{prop:lim_1}
    Sean $f:A  \subseteq  \Rn{n} \to \R$ una funci\'on y $\mathbf{x}_0$ un punto de acumulaci\'on de $A$. Las siguientes proposiciones son equivalentes:
 \begin{enumerate} 
  \item[a)] $\lim_{\mathbf{x} \to \mathbf{x}_0} f(\mathbf{x}) = 0$.
  \item[b)] $\lim_{\mathbf{x} \to \mathbf{x}_0} | f(\mathbf{x}) | = 0$.
 \end{enumerate}
\end{propertie} 
\begin{proof}   Veamos que $a) \Rightarrow b).$  Como 
  \[
   \lim_{\mathbf{x} \to \mathbf{x}_0} f(\mathbf{x}) = 0,
  \]
  por definici\'on de l\'imite, dado $\varepsilon > 0$ podemos tomar $\delta > 0$ tal que 
  \[
   \abs{f(\mathbf{x}) - 0} < \varepsilon \; \forall \mathbf{x} \in B^*_{\delta}(\mathbf{x}_0) \cap A.
  \]
  Para este valor de $\delta$ se cumple que
  \[
   \Big| \abs{f(\mathbf{x})} - 0 \Big| = \Big| \abs{f(\mathbf{x})} \Big| = \abs{f(\mathbf{x})} = \abs{f(\mathbf{x}) - 0} < \varepsilon \; 
   \forall \mathbf{x} \in B^*_{\delta}(\mathbf{x}_0) \cap A,
  \]
  de modo que, por definici\'on, 
  \[
   \lim_{\mathbf{x} \to \mathbf{x}_0} \abs{ f(\mathbf{x}) } = 0.
  \]
  
Veamos que  $b) \Rightarrow a).$  Como
  \[
   \lim_{\mathbf{x} \to \mathbf{x}_0} \abs{ f(\mathbf{x}) } = 0,
  \]
  por definici\'on de l\'imite, dado $\varepsilon > 0$ podemos tomar $\delta > 0$ tal que 
  \[
   \Big| \abs{f(\mathbf{x})} - 0 \Big| < \varepsilon \; \forall \mathbf{x} \in B^*_{\delta}(\mathbf{x}_0) \cap A.
  \]
  Para este valor de $\delta$ se cumple que
  \[
   \abs{f(\mathbf{x}) - 0} = \abs{f(\mathbf{x})} = \Big| \abs{f(\mathbf{x})} \Big| = \Big| \abs{f(x)} - 0 \Big| < \varepsilon \, 
   \forall \mathbf{x} \in B^*_{\delta}(\mathbf{x}_0) \cap A,
  \]
  de modo que, por definici\'on, 
  \[
   \lim_{\mathbf{x} \to \mathbf{x}_0} f(\mathbf{x}) = 0.
  \]
 \end{proof}



\begin{propertie} \label{prop:sandwich} %[Principio de intercalaci\'on]
Sean $f, g, h : A  \subseteq  \Rn{n} \to \R$ funciones y $\mathbf{x}_0$ un punto de acumulaci\'on de $A$. Supongamos que $\exists\: \delta_0 > 0$ tal que 
\[
 f(\mathbf{x}) \le g(\mathbf{x}) \le h(\mathbf{x}) \; \forall \mathbf{x} \in B_{\delta_0}(\mathbf{x}_0) \cap A.
\]
Si $\lim_{\mathbf{x} \to \mathbf{x}_0}f(\mathbf{x}) = \lim_{\mathbf{x} \to \mathbf{x}_0}h(\mathbf{x}) = L$, entonces
\[
 \lim_{\mathbf{x} \to \mathbf{x}_0}g(\mathbf{x}) = L.
\]
 
\end{propertie}

\begin{propertie} \label{prop:cero_x_acotada}
 Sean $f, g : A  \subseteq  \Rn{n} \to \R$ funciones y $\mathbf{x}_0$ un punto de acumulaci\'on de $A$. Si $\exists\: \delta_0 > 0$ tal que $f$ es acotada en $B_{\delta_0}(\mathbf{x}_0) \cap A$ (\textit{i.e.} $\exists\: M > 0 \text{ tal que } \abs{f(\mathbf{x})} \le M,\;\forall \mathbf{x} \in  B_{\delta_0}(\mathbf{x}_0) \cap A$) y $\lim_{\mathbf{x} \to \mathbf{x}_0}g(\mathbf{x}) = 0$, entonces
\[
 \lim_{\mathbf{x} \to \mathbf{x}_0}(f  g)(\mathbf{x}) = 0.
\]
\begin{proof} Supongamos que, para cierto $\delta_0 > 0$, $f$ es acotada en $B_{\delta_0}(\mathbf{x}_0) \cap A$, y sea $M > 0 \text{ tal que } \abs{f(\mathbf{x})} \le M, \; \forall \mathbf{x} \in  B_{\delta_0}(\mathbf{x}_0) \cap A$.
Como 
\[
\lim_{\mathbf{x} \to \mathbf{x}_0}g(\mathbf{x}) = 0, 
\]
dado $\varepsilon > 0 $, por definici\'on de l\'imite, podemos elegir $\delta_1 > 0$ tal que 
\[
 \abs{g(\mathbf{x})} < \frac{\varepsilon}{M},\;\forall \mathbf{x} \in  B^*_{\delta_1}(\mathbf{x}_0) \cap A.                                                                                               
\]
Sea $\delta = \min{ \left\{ \delta_0, \delta_1 \right\} }$. \\
Se cumple que
\[
 \abs{ (fg)(\mathbf{x}) - 0 } = \abs{ f(\mathbf{x}) \cdot g(\mathbf{x}) } = \abs{f(\mathbf{x})}\abs{g(\mathbf{x})} \le 
 M \abs{g(\mathbf{x})} < M \frac{\varepsilon}{M} = \varepsilon \quad \forall \mathbf{x} \in  B^*_{\delta}(\mathbf{x}_0) \cap A,
\]
Por lo tanto, por definici\'on de l\'imite,
\[
 \lim_{\mathbf{x} \to \mathbf{x}_0}(fg)(\mathbf{x}) = 0.
\]\end{proof}
\end{propertie}


\begin{propertie} \textbf{L\'imite de la composici\'on.} \label{prop:lim_comp} Sean $\mathbf{F}$ y $\mathbf{G}$ funciones 
\begin{align*}
 \mathbf{G} &:A  \subseteq  \Rn{n} \to B  \subseteq  \Rn{m} \\
 \mathbf{F} &:B  \subseteq  \Rn{m} \to \Rn{p},
\end{align*}
y puntos $\mathbf{a} \in A', \mathbf{b} \in B'$ tales que 
\[
 \mathbf{G}(\mathbf{x})\neq  \mathbf{b} \; \forall  \mathbf{x}\neq\mathbf{a},    \;\; \lim_{\mathbf{x} \to \mathbf{a}} \mathbf{G}(\mathbf{x}) = \mathbf{b} \quad \wedge \quad \lim_{\mathbf{u} \to \mathbf{b}} \mathbf{F}(\mathbf{u}) = \mathbf{L}.
\]
Entonces
\[
 \lim_{\mathbf{x} \to \mathbf{a}} \left( \mathbf{F} \circ \mathbf{G} \right)(\mathbf{x}) = \mathbf{L}.
\]
\end{propertie}
\begin{example}
 Veamos que 
 \[
  \lim_{(x,y) \to (0,0)} \frac{\sin(x^2 + y^2)}{x^2 + y^2} = 1.
 \]
 En efecto, sabemos que 
  \[
    \lim_{u \to 0} \frac{\sin u}{u} = 1,
  \]
  y que
  \[
    \lim_{(x,y) \to (0,0)} x^2 + y^2 = 0.
  \]  
  Por lo tanto, si definimos 
  \begin{gather*}
   f: \R \smallsetminus \{0\} \to \R \;|\;f(u) = \frac{\sin u}{u} \\
   G: \Rn{2} \smallsetminus \{(0,0)\} \to \R \;|\; G(x,y) = x^2 + y^2,
  \end{gather*}
  por la Propiedad \ref{prop:lim_comp} es
  \[
   \lim_{(x,y) \to (0,0)} \left( f \circ G \right)(x,y) = 
   \lim_{(x,y) \to (0,0)} \frac{\sin(x^2 + y^2)}{x^2 + y^2} = 1.
  \]
\end{example}

\begin{tarea}  Calcular, si existe,   $$\lim_{(x,y)\rightarrow (0,0)}  \dfrac{\sin^{2}(x^{\frac{1}{3}}y)   (e^{ x^{2}+ y^{2} }-1)}{\big( x^{2}+y^{2} \big)^{2}}$$
\end{tarea}


%Una consecuencia directa de la Propiedad \ref{prop:lim_comp} es el siguiente corolario: 
\begin{corollary} \textbf{L\'imites por curvas.} \label{cor:lim_curva}
  Sean $f: A  \subseteq  \Rn{2} \to \R, \, (x_0,y_0)$ un punto de acumulaci\'on de $A$ y $L \in \R$ tales que 
  \[
   \lim_{(x,y) \to (x_0,y_0)} f(x,y) = L.
  \]

  Entonces, para cualesquiera funciones continuas $g_1, \, g_2 : I  \subseteq  \R \to \R$ (donde $I$ es un intervalo) tales que, para alg\'un $t_0 \in I'$
  \[
    \lim_{t \to t_0} \left( g_1(t),g_2(t) \right) = \left( x_0 , y_0 \right)
  \]
$$ \left( g_1(t),g_2(t) \right) \neq (x_0, y_0) \;\mbox{ si }  t\neq t_0$$
  $$\left( g_1(t), g_2(t) \right) \in A \, \forall t \in \left( I \smallsetminus \{t_0\} \right),$$ se cumple que
  \[
   \lim_{t \to t_0} f\left( g_1(t),g_2(t) \right) = L.
  \]
\end{corollary}


\begin{tarea}Siendo  $$\lim_{(x,y)\rightarrow (0,0)} \dfrac{ 3x^{2} + 5xy^{2}+3y^{2}}{x^{2} + y^{2}} = l$$
\begin{itemize}
\item[1.] Usando  (\ref{cor:lim_curva})  ver que $l=3$  es un "buen" candidato a l\'imite. 
\item[2.] Probar que $l=3.$
\end{itemize}
\end{tarea}



\begin{example} Probar que $f$ no tiene l\'imite para $(x,y) \to (0,0)$ siendo $$ f: \left\lbrace (x,y) \in \Rn{2} : x \ne y \right\rbrace \to \R \: |\:  f(x,y) = \frac{x^2}{x - y}.$$
   \begin{enumerate} 
   \item  Sean $g_1(t) = t, \,  g_2(t) = t + t^2$.
  \[
   f(g_1(t),g_2(t)) = f(t,t + t^2) = \frac{t^2}{-t^2} \to -1 \text{ cuando } t \to 0,
  \]
  de modo que, \emph{si $f$ tiene l\'imite}, el l\'imite es -1, por \eqref{cor:lim_curva} y la unicidad del l\'imite.
  \item  Sean $g_1(t) = t, \, g_2(t) = t - t^2$.
  \[
   f(g_1(t),g_2(t)) = f(t,t + t^2) = \frac{t^2}{t^2} \to 1 \text{ cuando } t \to 0,
  \]
  de modo que, \emph{si $f$ tiene l\'imite}, el l\'imite es 1, por \eqref{cor:lim_curva} y la unicidad del l\'imite.
  \end{enumerate}
  Por lo tanto, por la unicidad del l\'imite \eqref{teo:unicidad_limite}, \emph{si $f$ tiene l\'imite $L$}, $L = 1$ y $ L= -1$, lo que es un absurdo.  Por lo cual $f$ no puede tener l\'imite cuando  $(x,y) \to (0,0)$  .
\end{example}



\begin{propertie}\textbf{L\'imites iterados.} \label{prop:lim_ite}
  Sean $f: A  \subseteq  \Rn{2} \to \R$, $(x_0,y_0)$ un punto de acumulaci\'on de $A$ y $L \in \R$ tales que
  \[
   \lim_{(x,y) \to (x_0,y_0)} f(x,y) = L.
  \]
  \begin{enumerate} %[I.]
   \item Si para cada $x \ne x_0$ existe el l\'imite
  \[
   \lim_{y \to y_0} f(x,y) = g(x)
  \]
  y adem\'as existe el l\'imite
  \[
   \lim_{x \to x_0} g(x) = \lim_{x \to x_0} \left( \lim_{y \to y_0} f(x,y) \right) = L_1,
  \]  
  entonces $L = L1$.
  \item Si para cada $y \ne y_0$ existe el l\'imite
  \[
   \lim_{x \to x_0} f(x,y) = h(y)
  \]
  y adem\'as existe el l\'imite
  \[
   \lim_{y \to y_0} h(y) = \lim_{y \to y_0} \left( \lim_{x \to x_0} f(x,y) \right) = L_2, 
  \]
  entonces $L = L_2$.
  \end{enumerate}
\end{propertie}
% \emph{Ejemplos:
% \begin{itemize}
%  \item existan los iterados, sean distintos
%  \item existan los iterados, sean iguales, pero no hay l\'imite
%  \item existan los iterados, sean iguales, probar por def que existe el l\'imite
%  \item ejemplo en el que exista el l\'imite doble pero no los iterados
% \end{itemize}
% }

\begin{propertie}\textbf{L\'imites por conjuntos.} \label{prop:lim_x_conj}
   Sea $\mathbf{F} : A  \subseteq  \Rn{n} \to \Rn{m}$ una funci\'on, con $A = A_1 \cup A_2$.\\
   Sea $\mathbf{F}_1$ la \emph{restricci\'on de $\mathbf{F}$ a $A_1$}:
 \begin{align*}
  \mathbf{F}_1 &: A_1  \subseteq  \Rn{n} \to \Rn{m} \\
  &\mathbf{F}_1(\mathbf{x}) = \mathbf{F}(\mathbf{x}).
 \end{align*}
Sea $\mathbf{F}_2$ la \emph{restricci\'on de $\mathbf{F}$ a $A_2$}:
 \begin{align*}
  \mathbf{F}_2 &: A_2  \subseteq  \Rn{n} \to \Rn{m} \\
  &\mathbf{F}_2(\mathbf{x}) = \mathbf{F}(\mathbf{x}).
 \end{align*}
 Sean $\mathbf{x}_0 \in A_1' \cap A_2'$ y $\mathbf{L} \in \Rn{m}$. Entonces
 \[
  \lim_{\mathbf{x} \to \mathbf{x}_0} \mathbf{F}(\mathbf{x}) = \mathbf{L} \iff 
  \lim_{\mathbf{x} \to \mathbf{x}_0} \mathbf{F}_1(\mathbf{x}) = \mathbf{L} \wedge \lim_{\mathbf{x} \to \mathbf{x}_0} \mathbf{F}_2(\mathbf{x}) = \mathbf{L}.
 \]
\end{propertie}


\begin{example}   Sea 
  \[
   f(x,y) = 
     \begin{dcases}
        \dfrac{e^x - 1}{x}  & \text{ si } x \ne 0  \\
         1                 & \text{ si } x = 0
     \end{dcases}
  \]
 Veamos que $\lim_{(x,y) \to (0,0)} f(x,y) = 1$.\\
 Sean $A_1 = \{ (x,y) \in \Rn{2} : x \ne 0\}, \, A_2 = \{ (x,y) \in \Rn{2} : x = 0\}$ y 
 \begin{align*}
  f_1 &: A_1 \to \R,\quad f_1(x,y) = \frac{e^x - 1}{x}; \\
  f_2 &: A_2 \to \R,\quad f_2(x,y) = 1.
  \end{align*}
 Como 
 \[
    \lim_{t \to 0} \frac{e^t - 1}{t} = 1
 \]
 y 
 \[
    \lim_{(x,y) \to (0,0)} x = 0,
 \]
 por la Propiedad \ref{prop:lim_comp} es 
 \[
    \lim_{(x,y) \to (0,0)} f_1(x,y) = 1.
 \]
Como adem\'as $\lim_{(x,y) \to (0,0)} f_2(x,y) = 1$ y $A_1 \cup A_2 = \Rn{2}$, por la Propiedad \ref{prop:lim_x_conj} es
 \[
  \lim_{(x,y) \to (0,0)} f(x,y) = 1.
 \]
\end{example}


\begin{tarea} Calcular, si existe,   $\lim_{(x,y) \to (0,0)} f(x,y)$, siendo  
  
 $$f(x,y) = 
     \begin{dcases}
        \dfrac{x^6 y}{x^6 + y^3}  & \text{ si }  y \ne -x^{2}  \\
         0               & \text{ si }   y  =  -x^{2}
     \end{dcases} $$
Sugerencia: Usar  (\ref{prop:lim_x_conj})  y la curva  $\alpha(t) = (t,-t^{2} + t^{4})$ 
\end{tarea}




\begin{definition}\textbf{L\'imite infinito.} \label{def:lim_inf}
  Sean $f : A  \subseteq  \Rn{n} \to \R$ una funci\'on y $\mathbf{x}_0$ un punto de acumulaci\'on de $A$. 
  \begin{enumerate} %[I.]
   \item Decimos que \emph{el l\'imite de $f$ para $\mathbf{x}$ tendiendo a $\mathbf{x}_0$ es $+\infty$} o que \emph{$f$ tiende a $+\infty$ cuando $\mathbf{x}$ tiende a $\mathbf{x}_0$} si
    \[
      \forall M > 0\;\;\exists\: \delta > 0\;|\;f(\mathbf{x}) > M, \; \forall 
	  \mathbf{x} \in B_{\delta}^*(\mathbf{x}_0) \cap A.
    \]
    En este caso, utilizamos la notaci\'on
    \[
      f(\mathbf{x}) \xrightarrow[\mathbf{x} \to \mathbf{x}_0]{} +\infty \quad \text{o} \quad \lim_{\mathbf{x} \to \mathbf{x}_0} f(\mathbf{x}) = +\infty.
    \]
   \item Decimos que \emph{el l\'imite de $f$ para $\mathbf{x}$ tendiendo a $\mathbf{x}_0$ es $-\infty$} o que \emph{$f$ tiende a $-\infty$ cuando $\mathbf{x}$ tiende a $\mathbf{x}_0$} si
    \[
      \forall M > 0\;\;\exists\: \delta > 0 \;|\; f(\mathbf{x}) < -M, \; \forall 
	  \mathbf{x} \in B_{\delta}^*(\mathbf{x}_0) \cap A.
    \]
    En este caso, utilizamos la notaci\'on
    \[
      f(\mathbf{x}) \xrightarrow[\mathbf{x} \to \mathbf{x}_0]{} -\infty \quad \text{o} \quad \lim_{\mathbf{x} \to \mathbf{x}_0} f(\mathbf{x}) = -\infty.
    \]
   \item Decimos que \emph{el l\'imite de $f$ para $\mathbf{x}$ tendiendo a $\mathbf{x}_0$ es $\infty$} o que \emph{$f$ tiende a $\infty$ cuando $\mathbf{x}$ tiende a $\mathbf{x}_0$} si
    \[
      \forall M > 0 \;\; \exists\: \delta > 0 \;|\; \abs{f(\mathbf{x})} > M.\; \forall \mathbf{x} \in B_{\delta}^*(\mathbf{x}_0) \cap A.
    \]
    En este caso, utilizamos la notaci\'on
    \[
      f(\mathbf{x}) \xrightarrow[\mathbf{x} \to \mathbf{x}_0]{} \infty \quad \text{o} \quad \lim_{\mathbf{x} \to \mathbf{x}_0} f(\mathbf{x}) = \infty.
    \]
  \end{enumerate}
\end{definition}

%\begin{definition}\textbf{L\'imite en el infinito.} \label{def:lim_x_inf}
%Sean $f : \Rn{n} \to \R$ una funci\'on y $L \in \R$. Decimos que
%la manera de presentar esta definici\'on no es consistente con las anteriores
%\[
 %f(\mathbf{x}) \xrightarrow[\norm{\mathbf{x}} \to +\infty]{} L \quad \text{o} \quad \lim_{\norm{\mathbf{x}} \to +\infty} f(\mathbf{x}) = L
%\]
%si 
%\[ 
 %\forall \varepsilon > 0 \;\;\exists\: \delta > 0 \;|\; f(\mathbf{x}) \in B_{\varepsilon}(L),\; \forall 
% \mathbf{x} \;|\; \norm{\mathbf{x}} > \delta.
%\]
%\end{definition}


\iffalse
\begin{propertie} \label{prop:lim_2}
 Sean $f:A  \subseteq  \Rn{m} \to \Rn{m}$, $g:A \to \R$ y $x_0$ un punto de acumulaci\'on de $A$, tales que 
 \[
  \lim_{x \to x_0} g(x) = L_g \quad \wedge \quad \lim_{x \to x_0} \left( (f + g)_{(x)} \right) = L.
 \]
 Entonces, existe el l\'imite $\lim_{x \to x_0} f(x)$ y, adem\'as, $\lim_{x \to x_0} f(x) = L - L_g$.
\end{propertie}
\fi
